\slajdid{03-s009}
\begin{frame}[t,label=03-s009]
\gusttekst
\begin{columns}[c,onlytextwidth]
\begin{column}{.47\textwidth}\centering
\begingroup
% Ujednačen mali simbol sa 8 pt oznakama; selekcioni priključci su levo.
% \DEdek{ime}{x}{y}; globalno DEbar/DEbubble menjaju izlaze.
\providecommand{\DEbar}{0}
\providecommand{\DEbubble}{0}
\newcommand{\DEdek}[3]{%
\begin{scope}[shift={(#2,#3)},font=\fontsize{8}{8.5}\selectfont]
\draw[DEvod] (0,0) rectangle (1.7,1.65);
\node[align=center,inner sep=1pt] at (.58,1.27) {Dekoder\\2/4};
\foreach \n/\y in {3/1.24,2/.96,1/.68,0/.40}{
 \ifnum\DEbar=1
  \node[anchor=east,inner sep=1pt] at (1.7,\y) {$\overline{\mathrm{Y\n}}$};
 \else
  \node[anchor=east,inner sep=1pt] at (1.7,\y) {Y\n};
 \fi
 \ifnum\DEbubble=1
  \draw[DEvod,fill=white] (1.75,\y) circle(.05);
  \draw[DEvod] (1.80,\y)--(1.95,\y);
 \else
  \draw[DEvod] (1.7,\y)--(1.95,\y);
 \fi
 \coordinate (#1-y\n) at (1.95,\y);
}
\foreach \lab/\yy in {S1/.96,S0/.68,CS/.40}{
 \node[anchor=west,inner sep=1pt] at (0,\yy) {\lab};
 \draw[DEvod] (-.2,\yy)--(0,\yy);
 \coordinate (#1-\lab) at (-.2,\yy);
}
\end{scope}}

\begin{tikzpicture}[x=1cm,y=1cm,font=\fontsize{8}{9}\selectfont]
\draw[DEvod] (-.12,-.45) rectangle (4.76,7.3);
\node[align=center] at (.75,6.85) {Dekoder\\4/16};
\DEdek{r}{0}{3}
\foreach \i/\y in {3/5.5,2/3.7,1/1.9,0/.1}{
 \DEdek{d\i}{3}{\y}
 \foreach \j in {3,2,1,0}{
  \pgfmathtruncatemacro{\izlaz}{4*\i+\j}
  \draw[DEvod] (d\i-y\j)--++(.15,0) node[right,inner sep=1pt] {Y\izlaz};
 }
 \draw[DEvod] (d\i-S1)--(2.45,0 |- d\i-S1);
 \draw[DEvod] (d\i-S0)--(2.68,0 |- d\i-S0);
}
\draw[DEvod] (r-y3)--(2.1,4.24)|-(d3-CS);
\draw[DEvod] (r-y2)--(2.22,3.96)|-(d2-CS);
\draw[DEvod] (r-y1)--(2.1,3.68)|-(d1-CS);
\draw[DEvod] (r-y0)--(2.22,3.4)|-(d0-CS);
\draw[DEvod] (2.45,6.46)--(2.45,-.16)--(-.55,-.16) node[left] {S1};
\draw[DEvod] (2.68,6.18)--(2.68,-.35)--(-.55,-.35) node[left] {S0};
\draw[DEvod] (r-S1)--(-.55,3.96) node[left] {S3};
\draw[DEvod] (r-S0)--(-.55,3.68) node[left] {S2};
\draw[DEvod] (r-CS)--(-.55,3.40) node[left] {CS};
\end{tikzpicture}
\endgroup

\end{column}
\begin{column}{.50\textwidth}
Piramidalna struktura.

U ovom slučaju nam je dovoljan jedan selekcioni signal na komponentama

Mogući su i drugačiji načini povezivanja internih selekcionih signala. Prikazani način dozvoljava da se na lak način identifikuju indeksi novih izlaza. Po pravilu u izlaznom nivou se koriste selekcioni signali „najniže težine, a onda u prethodnim nivoima sve veće težine. Uočiti da na ovaj način možemo realizovati i komponentu sa većim brojem izlaza jedino će se povećavati broj nivoa. Na primer komponenta \(6/2^6\) realizovana sa dekoderima 2/4 bi imala tri nivoa. Naziv piramidalna struktura je zbog oblika koju ova struktura ima kao i načina dekodovanja.
\end{column}
\end{columns}
\end{frame}
